\scanpage{25}
% archive sheet number: 24

\uline{NB} L'antipodisme peut être défini \inserted{intrinsèquement} indépendamment de tte condition de caractéristique, et quelle que soit la parité de $n$, dès que $\rho'$ est inversible, \inserted{grâce à (6), (7)}. \deleted{(Mais pour $n$ impair et} \inserted{Pour que $u$ soit} l'identité, il faut s. que $n$ impair et $S$ de car. $2$. Donc pour qu'il soit $\ne\mathrm{id}$ en chaque fibre, il f. et s. que $n$ pair ou $2$ inv. sur $S$.

\begin{center}\rule{1.1cm}{0.4pt}\\[-.6ex]\rule{.9cm}{0.4pt}\end{center}

Considérons $G\xrightarrow{\varphi_G}\operatorname{Aut}(X)$ satisfaisant les conditions habituelles (mais on ne donne pas pour l'\unsure{instt|instant} de réal$^{\mathrm n}$ géom. corr., i.e. pas de $s_0$\ldots). Donc on a $C\subset X$ permettant de construire (
\diagram{\begin{tikzcd}[ampersand replacement=\&]
0\arrow[r]\&V\arrow[r]\&E\arrow[r]\&\underline{\mathcal O}_S\arrow[r]\&0
\end{tikzcd}}
\margin{$X=\check{\mathbb P}(\unsure{E|\mathcal E})$}

et on a $L_i\subset V,\ L'_i\subset\check V_{\deleted{\unsure{S}}}$ $(0\leq i\leq n)$, \deleted{$L_i\subset\check V$} tels que $V=\bigoplus L_i$, \deleted{$\unsure{V'|\check V}=\bigoplus L'_i$}, $L_i\otimes L'_{i+1}\overset{\sim}{\longrightarrow}\underline{\mathcal O}_S$ pour $0\leq i\leq n-1$. \deleted{En \unsure{f|p} Considérons la} \inserted{Ainsi,} $G$ opère sur $V$ \deleted{\unsure{via}} de façon compatible avec son op sur $C\simeq\check{\mathrm P}(V)$. On a en \unsure{loc$^{\mathrm{ment}}$|laissant|outre} la donnée de $s_0$ (satisfaisant aux conditions habi%
